%% 第六章 结论与展望

\chapter{结论与展望}
\chaptermark{结论与展望}

\section{主要工作与结论}

本文面向载人登月任务需求，选取猎户座飞船为参考构型，开展了发射逃逸阶段的六自由度动力学建模、蒙特卡洛仿真与风险评估研究。主要工作与结论如下~\cite{刁伟鹤2025,Antonsen2023}：

（1）建立了发射逃逸六自由度动力学仿真模型。完成了坐标系定义与坐标变换、ISA大气分层模型、马赫数分段气动力/力矩模型、AM/ACM/JM多发动机联合推进模型、程序角PD姿控模型，以及六自由度全耦合运动微分方程组的推导，实现了Pad Abort和Max-Q Abort两种工况的数值仿真。

（2）建立了六类干扰项模型。包括风扰（水平风速、垂直风速和风向）、发动机安装偏差、推力幅值偏差、质量偏差、初始状态偏差，以及气动和大气模型偏差，共8个独立实验因素。所有干扰项的$3\sigma$对应工程极限偏差。

（3）完成了蒙特卡洛打靶仿真与统计分析。对两种工况分别开展了$L_{27}(3^8)$正交试验和500组拉丁超立方试验。零高度逃逸命中率96.20\%（命中半径50 m），最大动压逃逸命中率95.20\%（命中半径1000 m），在当前扰动口径下具备较高的命中精度。

（4）揭示了两工况的敏感性差异规律。零高度逃逸对风扰最为敏感（风向极差13.00 m），物理原因在于地面低速飞行阶段气动阻尼弱、姿控系统需直接对抗风扰。最大动压逃逸对气动和大气模型偏差最为敏感（极差276.30 m），物理原因在于初始高动压条件使气动参数不确定性放大为显著的弹道偏差。

（5）完成了风险评估。建立以超限概率为核心的分级评估体系。两种工况的命中风险和姿态风险均为低风险等级（超限概率$<4\%$，角速度变化$<0.003$ rad/s）。最大动压逃逸的尾部风险值得关注（最大脱靶量1414.684 m）。

\section{不足与展望}

本文工作仍存在以下不足，可作为后续研究的方向：

\begin{enumerate}
    \item 模型参数精度。本文气动系数基于马赫数三分段拟合，实际跨声速区的气动特性更为复杂。若获取更精细的风洞试验数据或CFD计算结果，可大幅提升仿真精度。一些数据来自于估算，存在误差。

    \item 干扰项相关性。本文假设8个实验因素相互独立，实际工程中部分干扰项（如推力偏差与质量偏差）可能存在相关性。引入协方差矩阵可更准确地反映真实扰动分布。

    \item 评估工况扩展。本文仅覆盖了发射逃逸段的Pad Abort和Max-Q Abort两种工况。后续可将仿真范围扩展到高空逃逸、抛塔后整船逃逸等阶段，形成完整的全任务链仿真与分析框架。

    \item 控制律优化。本文采用程序角比例微分（PD）控制律，增益为固定值。未来可引入自适应控制或模型预测控制方法，进一步提高系统对扰动的鲁棒性。

    \item 蒙特卡洛效率。拉丁超立方试验500组的统计量已基本满足工程分析需求，但在置信区间较窄的尾部风险分析中，可考虑增加样本量或采用重要性抽样等方法提高效率。
\end{enumerate}
